\documentclass[11pt,a4paper]{moderncv}
\usepackage{times}
\usepackage{xcolor}
\usepackage{hyperref}
\hypersetup{
	colorlinks=true,
	linkcolor=teal,
	citecolor=teal,
	filecolor=teal,      
	urlcolor=teal,
	pdftitle={Overleaf Example},
	pdfpagemode=FullScreen,
}
\usepackage[a4paper, total={7in, 10in}]{geometry}
\moderncvstyle{classic}   % style
\moderncvcolor{blue}      % couleur

\usepackage[utf8]{inputenc}

\name{Feriel}{Bouhadjera}
\title{Ma\^itresse de conférences}
\address{CNAM, 2 rue Conté, 75003 Paris}
\phone{01 40 27 23 73}
\email{prenom.nom@lecnam.net}

\begin{document}

\makecvtitle

\section{Cursus}
\cventry{2022--Auj}{Ma\^itresse de conférences}{\href{https://www.cnam.fr/}{Conservatoire National des Arts et Métiers}}{Laboratoire \href{https://cedric.cnam.fr/}{CEDRIC}}{Paris}{}
\cventry{2021--2022}{Post-doctorante en statistique}{\href{https://www.inrae.fr/}{Institut National de Recherche pour l'Agriculture, l'Alimentation et l'Environnement}}{Laboratoire de \href{https://mistea.montpellier.hub.inrae.fr/}{Mathématiques, Informatique et Statistique pour l’Environnement et l’Agronomie}}{Montpellier}{}
\cventry{2020--2021}{Attachée Temporaire d'Enseignement et de Recherche}{\href{https://www.univ-lille.fr/}{Université de Lille}}{Laboratoire \href{https://math.univ-lille.fr/}{Paul Painlevé}}{Lille}{}

\cventry{2018--2020}{Doctorante en statistique mathématique}{\href{https://www.univ-littoral.fr/}{Université du Littoral Côte d'Opale}}{Laboratoire de \href{https://lmpa.univ-littoral.fr/}{Mathématiques pures et appliquées}}{Calais, France}{}
\cvitem{}{$\triangleright$ Titre : Estimation non paramétrique de la fonction de régression pour des données censurées : méthodes locale linéaire et erreur relative.}
\cvitem{}{$\triangleright$ Jury :  \textrm{Célestin Kokonendji} (PR, Université Bourgogne Franche-Comté) [examinateur, président], \textrm{Zohra Guessoum} (PR, Université des Sciences et de la Technologie Houari Boumédiène) [rapporteur], \textrm{Mustapha Rachdi} (PR, Université de Grenoble Alpes) [rapporteur], \textrm{Hacène Boutabia} (PR, Université Badji Mokhtar Annaba) [examinateur] -- \textrm{Ouafae Benrabah} (MCF, Université du Littoral C\^ote d'Opale) [invitée], \textrm{Elias Ould Sa\"id} (PR, Université du Littoral C\^ote d'Opale) [directeur], \textrm{Mohamed Riad Remita} (PR, Université Badji Mokhtar Annaba) [directeur].}
\cventry{2016--2017}{Doctorante en 1 ère année, spécialité Actuariat}{Université Badji Mokhtar Annaba, Laboratoire de probabilités et Statistique}{Annaba, Algérie}{}{}
\cvitem{}{$\triangleright$ Titre : Évaluation des options exotiques dans un marché généré par des agents
interactifs.}
\cventry{2014--2016}{Master en probabilités et statistique}{Université Badji Mokhtar Annaba}{Algérie}{}{}
\cvitem{}{$\triangleright$ Mention Très bien}
\cventry{2011--2014}{Licence en probabilités et statistique}{Université Badji Mokhtar Annaba}{Algérie}{}{}
\cvitem{}{$\triangleright$ Mention Très bien}
%----------------------------
\section{Vacations}
\cvitem{2022}{Vacataire à l'Institut Agro, Montpellier.}
\cvitem{2020}{Vacataire à l'Institut Supérieur de Gestion, Lille.}
\cvitem{2019--2020}{Vacataire à l'Université du Littoral C\^ote d'Opale, Calais.}
\cvitem{2016--2017}{Vacataire à l'Université Badji Mokhtar Annaba, Algérie.}
%----------------------------------
\section{Productions scientifiques}
\subsection{\textcolor{orange}{Publications}}
\cvitem{1.}{Bouhadjera F., Lemdani M. \& Ould Sa\"id E. (2023). Strong uniform consistency of the local linear relative error regression estimator under left truncation. \emph{Statistical Papers, 64(2), 421--447.} \url{https://doi.org/10.1007/s00362-022-01325-9}} 
\cvitem{2.}{Bouhadjera F. \& Ould Sa\"id E. (2023). Relative error prediction: Strong uniform consistency for censoring time series model. \emph{Communications in Statistics : Theory and Methods, 52(11), 3709--3729.} \url{https://doi.org/10.1080/03610926.2021.1979584}} 
\cvitem{3.}{Bouhadjera F., Ould Sa\"id E. \& Remita. M.R. (2022). Strong consistency of the nonparametric local linear regression estimation under censorship model. \emph{Communications in Statistics : Theory and Methods, 51(20), 7056--7072.}\url{https://doi.org/10.1080/03610926.2020.1870142}}
\cvitem{4.}{ Benrabah O., Bouhadjera F. \& Ould Sa\"id E. (2022). Local linear estimation of the regression function for twice censored data. \emph{Statistical Papers, 63, 486--514.} \url{https://doi.org/10.1007/s00362-021-01240-5}}
\cvitem{5.}{Bouhadjera F. \& Ould Sa\"id E. (2022). Nonparametric local linear estimation of the relative error regression function for censorship model. \emph{Opuscula Mathematica, 42(6), 805--832.} \url{https://doi.org/10.7494/OpMath.2022.42.6.805}}
\cvitem{6.}{Bouhadjera F. \& Ould Sa\"id E. (2021). Asymptotic normality of the relative error regression function estimator for censored time series data. \emph{Dependence Modeling, 9(1), 156--179.} \url{https://doi.org/10.1515/demo-2021-0107}}
\cvitem{7.}{Bouhadjera F. \& Ould Sa\"id E. (2021). Nonparametric local linear estimation of the relative error regression function for twice censored data. \emph{Statistics \& Probability Letters, 178, 109185.} \url{https://doi.org/10.1016/j.spl.2021.109185}} 
%------------------------------------
\subsection{\textcolor{orange}{Thèse}}
\cvitem{1.}{Bouhadjera, F. (2020). \emph{Estimation non paramétrique de la fonction de régression pour des données censurées : méthodes locale linéaire et erreur relative. \url{https://theses.hal.science/tel-03134914v1/file/these_Bouhafjera_Feriel.pdf}}}
%------------------------------------
\subsection{\textcolor{orange}{Autres travaux}}
\cvitem{1.}{Bouhadjera F., Ould Sa\"id E. \& Remita. M.R. (2019). Nonparametric relative error estimation of the
regression function for censored data. \emph{Arxiv}. \url{https://arxiv.org/pdf/1901.09555}}
%-------------------------------------

\subsection{\textcolor{orange}{Conférences Internationales}}
\cvitem{1.}{ Diallo, B., Niang, N., Audigier, V. \& Bouhadjera, F. (2026). Hybrid data clustering with derivative information. \emph{Advanced Machine Learning and Data Science}.}
\cvitem{2.}{Bouhadjera, F. \& Ould-Sa\"id  E. (2019). Nonparametric kernel estimation for survival right censored data. \emph{International Conference on Mathematical Modeling and Applications}. Rabat, Maroc.}
\cvitem{3.}{Bouhadjera, F. \& Remita M.R. (2017). Évaluation des options exotiques dans un marché généré par des agents interactifs. \emph{International conference on the evolution of the contemporary mathematics and its impact on science and technology, Université de Frères Mentouri, Constantine, Algérie.}}
%------------------------------------
\subsection{\textcolor{orange}{Conférences Nationales et Workshops}}
\cvitem{1.}{Diallo, B., Niang, N., Audigier, V., \& Bouhadjera, F. (2026). Comparaison de méthodes de classification de données fonctionnelles. \emph{Revue des Nouvelles Technologies de l'Information, Extraction et Gestion des Connaissances, RNTI-E-42: 121--132.} \url{https://hal.science/hal-05488620}}
\cvitem{2.}{ Diallo, B., Niang, N.,  Bouhadjera, F. \& Audigier, V. (2026). Classification de données fonctionnelles et vectorielles. \emph{57èmes Journées de la Statistique.} \url{https://hal.science/hal-05589708}}
\cvitem{3.}{Bouhadjera, F. (2024). Régression sur données censurées avec indicateurs manquants. \emph{55èmes Journées de la Statistique.} \url{https://hal.science/hal-05010361}}
\cvitem{4.}{Bouhadjera, F., Baragatti, M., Hilgert, N., Smits, N. \& Grollemund, P-M. (2022). Bayesian functional linear regression estimation. Extension to scalar and categorical covariates. \emph{3rd LmB Conference on Multivariate Statistical Models: Count and Semi-Continuous.}}
\cvitem{5.}{Bouhadjera, F., Baragatti, M., Hilgert, N., Smits, N. \& Grollemund, P-M. (2022). Bayesian functional linear regression estimation. Extension to scalar and categorical covariates. \emph{53èmes Journées de la Statistique.} \url{https://hal.inrae.fr/hal-04189804}}
\cvitem{6.}{Bouhadjera, F., Baragatti, M., Hilgert, N., Smits, N. \& Grollemund, P-M. (2021). Régression linéaire bayésienne sur données fonctionnelles. \emph{Consortuim Phenodyn. INRAE, Montpellier.}}
\cvitem{7.}{Bouhadjera, F. (2019). Propriétés asymptotiques de l'estimateur à noyaux en présence de censure. \emph{13ème rencontre doctorale en mathématiques du Nord-Pas-de-Calais}, Lens, France.}
\cvitem{8.}{Bouhadjera, F. (2019). Estimation non paramétrique de la fonction de régression relative en présence de censure. \emph{Journées des jeunes chercheurs du laboratoire de probabilités et statistique}, Annaba, Algérie.}
\cvitem{9.}{Bouhadjera, F. (2017). Estimation de la fonction de régression quantile via les copules. \emph{1ère édition des doctoriales nationale en mathématiques, Université des Frères Mentouri, Constantine, Algérie.}}
%------------------------------------
\subsection{\textcolor{orange}{Présentations dans des séminaires}}
\cvitem{1.}{Bouhadjera, F. (Juin, 2023). Modèles de régression pour données de survie. \emph{Séminaire du Laboratoire Paul Painlevé, Lille.} }
\cvitem{2.}{Bouhadjera, F. (Mai, 2023). Modèles de régression pour données de survie. \emph{Séminaire du Laboratoire de Mathématiques Appliquées de Compiègne.} }
\cvitem{3.}{Bouhadjera, F. (Octobre, 2022). Modèles de régression pour données de survie. \emph{Séminaire du Laboratoire CEDRIC, CNAM.} \url{https://maths.cnam.fr/IMG/pdf/expose_axe_cnam_bouhadjera.pdf}}
\cvitem{4.}{Bouhadjera, F. (Octobre, 2021). Régression non paramétrique pour des données incomplètes. \emph{Séminaire de l'Institut Montpelliérain Alexander Grothendieck, Université de Montpellier.}}
\cvitem{5.}{Bouhadjera, F. (Mai, 2021). Estimation non paramétrique locale linéaire de la fonction de régression relative pour des
données incomplètes. \emph{Séminaire du Laboratoire de Mathématiques et leurs Applications de Valenciennes.}}
\cvitem{6.}{Bouhadjera, F. (Mars, 2021). Propriétés asymptotiques de l'estimateur non paramétrique local linéaire de la fonction de régression pour des données censurées. \emph{Séminaire du Laboratoire Paul Painlevé, Lille.} }
\cvitem{7.}{Bouhadjera, F. (Mars, 2021). Estimation non paramétrique locale linéaire de la fonction de régression relative pour des données censurées. \emph{Séminaire de Statistique et Optimisation, Université de Toulouse.} }
\cvitem{8.}{Bouhadjera, F. (Février, 2021). Convergence uniforme de l'estimateur de la régression relative pour un modèle de censure aléatoire à droite. \emph{Séminaire du Laboratoire de Probabilités, Statistique et Applications, Université de Poitiers.} }
\cvitem{9.}{Bouhadjera, F. (Février, 2021). Estimation non paramétrique de la fonction de régression en présence de censure. \emph{Séminaire du Laboratoire de Statistique, Université de Strasbourg.} }
\cvitem{10.}{Bouhadjera, F. (Janvier, 2021). Estimation non paramétrique linéaire locale de la fonction de régression relative pour un modèle de double censure. \emph{Séminaire du Laboratoire de Probabilités, Statistique et Application, Université Houari Boumédiène, Alger, Algérie.}}
\cvitem{11.}{Bouhadjera, F. (Octobre, 2019). Estimation de la fonction de régression à noyaux pour des données censurées. }
\cvitem{13.}{Bouhadjera, F. (Décembre, 2018). Estimation non paramétrique de la fonction de régression relative en présence de censure. \emph{Séminaire du Laboratoire de Probabilités, Statistique et Application, Université Houari Boumédiène, Alger, Algérie.}}
\cvitem{13.}{Bouhadjera, F. (Novembre, 2018). Estimation non paramétrique de la fonction de régression relative en présence de censure. \emph{Séminaire du laboratoire de Mathématiques Pures et Appliquées}, Université du Littoral C\^ote d'Opale, Calais, France.}
%------------------------------------
\newpage
\section{Enseignements}
\cvitem{2025--Auj}{\textbf{Apprentissage statistique en grande dimension.} \textit{Formation en hybride}, CNAM, Paris.}
\cvitem{}{$\triangleright$ Formation continue, niveau M2}
\cvitem{}{$\triangleright$ Méthodes non-supervisées (ACP fonctionnelle, clustering fonctionnel) et méthodes de régression pour données fonctionnelles. (\href{https://www.cnam-paris.fr/choisir-ma-formation/par-discipline/apprentissage-statistique-en-grande-dimension-1551348.kjsp}{Lien vers le programme complet})}
%------------------------
\cvitem{2025--Auj}{\textbf{Introduction aux données fonctionnelles.} \textit{Formation à distance}, CNAM, Paris.}
\cvitem{}{$\triangleright$ Formation continue, niveau M1}
\cvitem{}{$\triangleright$ Pré-traitement de données fonctionnelles (représentation dans une base de fonctions, lissage), analyse descriptive (moyenne, variance, covariance fonctionnelles), outils de visualisation (Rainbowplot, Bagplot, HDRplot), analyse exploratoire (ACP fonctionnelle) et clustering fonctionnel. \href{https://cnam-idf.fr/unites/introduction-aux-donnees-fonctionnelles-sta117}{(Lien vers le programme complet)}}
%------------------------
\cvitem{2025--Auj}{\textbf{Modélisation et prévision des séries chronologiques.} \textit{Formation à distance}, CNAM, Paris.}
\cvitem{}{$\triangleright$ Formation continue, niveau M1}
\cvitem{}{$\triangleright$ Décomposition des séries chronologiques, lissages exponentiels (simple, double, Holt-Winters), modélisation des séries stationnaires et extension aux séries non stationnaires. \href{https://www.cnam-paris.fr/choisir-ma-formation/catalogue-de-la-foad/modelisation-et-prevision-des-series-chronologiques-1595208.kjsp?RH=1404475742767}{(Lien vers le programme complet)}}
%--------------------------
\cvitem{2023--2025}{\textbf{Biostatistique.} \textit{Formation en présentiel}, Pasteur Cnam, Paris.}
\cvitem{}{$\triangleright$ Mastère}
\cvitem{}{$\triangleright$ Probabilités élémentaires, variables aléatoires, estimation ponctuelle et par intervalle, tests statistiques.}
%------------------------
\cvitem{2022--Auj}{\textbf{Calcul des probabilités.} \textit{Formation en présentiel}, CNAM, Paris.}
\cvitem{}{$\triangleright$ Formation continue, niveau L2}
\cvitem{}{$\triangleright$ Notions de probabilités, variables aléatoires, lois usuelles, fonctions génératrice et caractéristique, lois empiriques et comportement asymptotique (convergence, lois des grands nombres, TCL). \href{https://www.cnam.fr/formation/sciences/mathematiques/calcul-des-probabilites-2}{(Lien vers le programme complet)}}
%------------------------
\cvitem{2022--Auj}{\textbf{Statistique mathématique.} \textit{Formation à distance}, CNAM, Paris.}
\cvitem{}{$\triangleright$ Formation continue, niveau L3}
\cvitem{}{$\triangleright$ Estimation paramétrique : méthode du maximum de vraisemblance, méthode des moments, intervalles de confiance. Estimation non paramétrique : méthodes à noyau, $k$-plus proches voisins, splines. Tests statistiques (khi-deux, kolmogorov, Wilcoxon.) \href{https://www.cnam-paris.fr/choisir-ma-formation/par-discipline/statistique-mathematique-1487936.kjsp}{(Lien vers le programme complet)}}
%--------------------------
\cvitem{2022--Auj}{\textbf{Analyse multivariée approfondie.} \textit{Formation à distance et présentiel}, CNAM, Paris.}
\cvitem{}{$\triangleright$ Formation continue, niveau M2}
\cvitem{}{$\triangleright$ Méthodes d'estimation non paramétriques de la densité et de la régression. \href{https://www.cnam-paris.fr/choisir-ma-formation/par-discipline/analyse-multivariee-approfondie-1487944.kjsp}{(Lien vers le programme complet)}}
%--------------------------
\cvitem{2021--2022}{\textbf{Statistique avec R.} \textit{Formation en présentiel}, Institut Agro, Montpellier.}
\cvitem{}{$\triangleright$ Ingénieur agronome, niveau L3}
\cvitem{}{$\triangleright$ Analyse statistique des données avec R Markdown (Tests statistiques, Modélisation et sélection de variables)}
%--------------------------
\cvitem{2020--2021}{\textbf{Mathématiques fondamentales.} \textit{Formation en présentiel}, Université de Lille, Lille.}
\cvitem{}{$\triangleright$ Niveau L1}
\cvitem{}{$\triangleright$ Limite, continuité, dérivabilité, intégration, équation différentielle,
espace vectoriel, application linéaire et calcul matriciel, développement limité}
%--------------------------
\cvitem{2020}{\textbf{Mathématiques financières.} \textit{Formation en présentiel}, Institut Supérieur de Gestion, Lille.}
\cvitem{}{$\triangleright$ Niveaux M1 et M2}
\cvitem{}{$\triangleright$ Valorisation d’options en finance de marché, modèle de Black and
Scholes, les options réelles et choix d’investissement, VAN augmenté}
%--------------------------
\cvitem{2019--2020}{\textbf{Mathématiques générales et statistiques.} \textit{Formation en présentiel}, Université du Littoral C\^ote d'Opale, Boulogne-sur-Mer.}
\cvitem{}{$\triangleright$ Niveau L1 en Économie et Gestion}
\cvitem{}{$\triangleright$ Matrice, systèmes d’équations linéaires, séries statistiques à une et
deux variables, séries chronologiques, mesure de la croissance}

%--------------------------
\cvitem{2016--2017}{\textbf{Biostatistique.} \textit{Formation en présentiel}, Université Badji Mokhtar Annaba, Algérie.}
\cvitem{}{$\triangleright$ Niveau L2 en Biologie}
\cvitem{}{$\triangleright$ Statistique descriptive, régression linéaire, tests (paramétrique / non
paramétrique)}
%--------------------------
\cvitem{2016--2017}{\textbf{Méthodes statistiques.} \textit{Formation en présentiel}, Université Badji Mokhtar Annaba, Algérie.}
\cvitem{}{$\triangleright$ Niveau M2 en Métallurgie}
\cvitem{}{$\triangleright$ Échantillonnage et estimation, régression linéaire, tests usuels et
analyse de variance}
%--------------------------
\cvitem{2016--2017}{\textbf{Mathématiques pour la physique.} \textit{Formation en présentiel}, Université Badji Mokhtar Annaba, Algérie.}
\cvitem{}{$\triangleright$ Niveau M1 en Physique}
\cvitem{}{$\triangleright$ EDP, dérivée et formule de Taylor, série de Fourier, transformé de
Fourier}

%---------------------------------
\section{Expériences professionnelles}

\subsection{\textcolor{orange}{Relecture d'articles}}
\textit{Journal of Computational and Applied Mathematics} et \textit{Journal KYBERNETIKA.}




%-----------------------
\subsection{\textcolor{orange}{Organisation de congrès internationaux}}
\cvitem{2025}{\textbf{Workshop on Multiblock Data Analysis and Related Methods.} CNAM, Paris, France. Membre du comité d’organisation}
\cvitem{2023}{\textbf{Symbolic Data Analysis Workshop.} CNAM, Paris, France. Membre du comité d’organisation}
%-----------------------
\subsection{\textcolor{orange}{Organisation de congrès nationaux}}
\cvitem{2024}{Demi-journée autour de l’eXplainable Artificial Intelligence (XAI). CNAM, Paris.}
%---------------------
\subsection{\textcolor{orange}{Implication dans des instances administratives}}
\cvitem{2024-Auj}{\textbf{Membre du conseil} de l'équipe pédagogique nationale de Mathématiques et Statistique. CNAM, Paris, France.}
\cvitem{2020-2021}{\textbf{Représentante des doctorants} au conseil du laboratoire de Mathématiques Pures et Appliquées. Université du Littoral C\^ote d'Opale, Calais, France.}

%------------------------------
\subsection{\textcolor{orange}{Participation à des comités de sélection}}
\cvitem{2025}{\textbf{MCF en Section 27}. IUT Toulouse Auch Castres, France.}


%-------------------------------
\subsection{\textcolor{orange}{Membre de comité de perfectionnement de dipl\^omes}}
\cvitem{2024--2025}{Master Science des données, CNAM}
\cvitem{2024--2025}{Licence Science des données, CNAM}

%----------------------------------
\subsection{\textcolor{orange}{Participation à des réseaux de recherche}}
\cvitem{2022}{DigitAgors, consortium sur les avancées scientifiques et pédagogiques en Agriculture et Numérique. Sète.}
\cvitem{2021--2022}{Groupe de travail Python/R du laboratoire MISTEA, INRAE, Montpellier.}
\cvitem{2021--2022}{Consortium Phenodyn du laboratoire MISTEA, INRAE, Montpellier.}
%-------------------------------
\subsection{\textcolor{orange}{Autres responsabilités}}
\cvitem{2026}{Modératrice pour la journées des doctorants de 2 ème année de l'école doctorale \href{https://adum.fr/script/formations.pl?mod=23111&menu_transparent=oui&site=smi}{SMI}.}
\cvitem{2025--Auj}{Co-organisation du séminaire de statistique de l'équipe MSDMA, CNAM, Paris}
\cvitem{2023}{Participation aux journées jeunes statisticiens et enseignement à l'IHP, Paris}

%-------------------------------
\section{Encadrements}
\subsection{\textcolor{orange}{Doctorants}}
\cvitem{2025--2027}{\textbf{Boubacar Diallo}, \textit{Classification non-supervisée pour l’identification de
paysages acoustiques homogènes.}}
\cvitem{}{$\triangleright$ Taux d'encadrement $35\%$}
\cvitem{}{$\triangleright$ Co-encadré avec Ndèye Niang et Vincent Audigier}
%--------------------------
\subsection{\textcolor{orange}{Stagiaires}}
\cvitem{2026}{\textbf{Sabin Bohissou.} M2. \textit{Classification supervisée de données spatio-temporelles avec
application aux capteurs tactiles.}}
\cvitem{2026}{\textbf{Roger Kouakou.} M2. \textit{Méthodes statistiques et de data-mining appliquées au
dispositif de jeu vidéo Tactile Game.}}
\cvitem{2026}{\textbf{Pierre Chambet.} M2. \textit{Classification non-supervisée de données hybrides.}}
\cvitem{2026}{\textbf{Lancelot Dallain, Tanguy Cesar.} M2. \textit{Classification non-supervisée de
données fonctionnelles avec application en océanographie.}}
\cvitem{2025}{\textbf{Hugo Nguyen}. M2. \textit{Classification non supervisée de paysages acoustiques.}}
\cvitem{2024}{\textbf{Marcellin Djohoundjo}. M2. \textit{Prédiction non paramétrique pour la géolocalisation d’objets.}}
\cvitem{2024}{\textbf{Djafarou Oulare}. M2. \textit{Classification non supervisée pour l’identification de paysages acoustiques homogènes.}}
\cvitem{2023}{\textbf{Koffi Amezouwui}. M2. \textit{Imputation multiple par régression non paramétrique.}}
\end{document}
